% Generated draft: requires mathematical and source review.
\documentclass[11pt]{article}
\usepackage[margin=25mm]{geometry}
\usepackage{fontspec}
\setmainfont{Latin Modern Roman}
\usepackage{amsmath,amssymb,mathtools,unicode-math}
\setmathfont{Latin Modern Math}
\usepackage{enumitem,needspace,xurl,hyperref}
\setlength{\parindent}{0pt}
\setlength{\parskip}{5pt}
\setlength{\emergencystretch}{3em}
\newcommand{\N}{\mathbb{N}}
\newcommand{\Z}{\mathbb{Z}}
\newcommand{\Q}{\mathbb{Q}}
\newcommand{\R}{\mathbb{R}}
\newcommand{\C}{\mathbb{C}}
\newcommand{\F}{\mathbb{F}}
\newcommand{\A}{\mathbb{A}}
\newcommand{\PP}{\mathbb{P}}
\newcommand{\E}{\mathbb{E}}
\newcommand{\eps}{\varepsilon}
\newcommand{\dd}{\,\mathrm{d}}
\begin{document}
\section*{UnsolvedMath ID: 30005596}
\textbf{Source problem code:} catalogue entry
\subsection{Definitions and mathematical statement}
\textbf{Supplied research question.} Determine sharp relationships between the lowest magnetic Neumann eigenvalue and the geometry of a planar domain, for constant magnetic fields and for Aharonov--Bohm potentials. The supplied record does not specify a particular inequality, admissible domain class, geometric constraint, or extremal problem. Consequently, it determines a research program rather than a single fully quantified conjecture.

\textbf{Regular Neumann framework.} To explain the spectral quantity in a standard regular setting, suppose conditionally that \(\Omega\subset\R^2\) is a bounded Lipschitz domain and that \(A\in C^1(\overline\Omega;\R^2)\) is a real magnetic vector potential. These regularity and boundedness hypotheses provide a reference definition; they are not supplied hypotheses of the source question. The magnetic Neumann Laplacian is the self-adjoint operator associated with the quadratic form
\[
q_{A,\Omega}[u]=\int_\Omega\lvert(-i\nabla-A)u\rvert^2\,\dd x,
\qquad \mathcal D(q_{A,\Omega})=H^1(\Omega;\C).
\]
Its lowest eigenvalue is
\[
\lambda_1^N(\Omega,A)=
\inf_{u\in H^1(\Omega;\C)\setminus\{0\}}
\frac{\displaystyle\int_\Omega\lvert(-i\nabla-A)u\rvert^2\,\dd x}
{\displaystyle\int_\Omega\lvert u\rvert^2\,\dd x}.
\]
When a classical boundary condition is available, the natural magnetic Neumann condition is
\[
\nu\cdot(-i\nabla-A)u=0\quad\text{on }\partial\Omega,
\]
where \(\nu\) is the outward unit normal.

A constant magnetic field of strength \(B\in\R\) has the representative potential
\[
A_B(x)=\frac B2(-x_2,x_1),
\qquad \partial_1(A_B)_2-\partial_2(A_B)_1=B.
\]
A standard single-pole Aharonov--Bohm potential, with pole \(p=(p_1,p_2)\in\R^2\) and circulation parameter \(\alpha\in\R\), is
\[
A_{\alpha,p}(x)=\alpha\frac{(-(x_2-p_2),\,x_1-p_1)}{|x-p|^2},
\qquad x\ne p.
\]
Its circulation around a positively oriented circle enclosing \(p\) is \(2\pi\alpha\). This formula explains the potential type; the record does not specify whether one or several poles are intended, their allowed locations, or their circulation parameters. If a pole meets \(\overline\Omega\), the regular form definition above cannot be applied without specifying the form domain and the treatment of the singularity. The record supplies no such specification.

\textbf{Unresolved formulation.} In each intended magnetic setting, the requested task is to identify sharp geometric estimates for the corresponding lowest Neumann eigenvalue. The domain hypotheses, permitted topology, magnetic parameters and their quantifiers, geometric quantities or constraints, direction of the estimates, and meaning of sharpness remain unspecified. No particular bound or extremizer is asserted here.
\subsection{Short English statement}
The supplied question asks how sharply planar-domain geometry controls the lowest magnetic Neumann eigenvalue for constant fields and Aharonov--Bohm potentials. It does not identify a specific inequality or optimization problem. The domain class, magnetic parameters, and treatment of singular poles require clarification before a fully quantified conjecture can be stated.
\subsection{Sources}
\begin{itemize}
\item Supplied local UnsolvedMath registry record: permanent ID \(30005596\), title ``Geometric Bounds for Magnetic-Laplacian Ground States,'' exact source problem code \texttt{catalogue entry}. Project registry paths: \texttt{data/unsolvedmath/problems.json} and \texttt{catalog/problems.json}. This draft uses the supplied complete record; the registry files were not independently inspected.
\item Supplied source citation: \emph{Geometric Spectral Theory (2024)}, Oberwolfach Reports, DOI \texttt{10.4171/owr/2023/36}, pp.~2017--2106. \url{https://doi.org/10.4171/owr/2023/36}. The record separately lists the proposed year as 2023.
\item Attribution in the supplied original excerpt: Bruno Colbois, ``Geometric aspects of the ground state of magnetic Laplacians on domains of the plane,'' joint work with C.~L\'ena, L.~Provenzano and A.~Savo. The excerpt poses the relation between the first eigenvalue and domain geometry as a natural question. The registry's proposer field is empty, so this attribution is not taken to establish a separately identified proposer.
\item Supplied literature source: ``Oberwolfach Report 36,'' \url{https://doi.org/10.4171/owr/2023/36}. Literature checked at: \texttt{2026-08-22}. This is supplied dated curation, not an independently verified assessment. It labels the problem open and reports that no general resolution was located; the background explicitly requests independent verification of time-sensitive status and cited claims.
\item The supplied record labels the statement \texttt{corrected\_verified} and says it was reconstructed and checked against the official report. Those are source-provided verification claims; neither the report nor those claims were independently checked for this draft.
\end{itemize}
\end{document}
